% answers_short.tex —— 1980 年·全国卷（文科） · 简短答案速查
% 由 tools/make_short_answers.py 自 answers.tex 的 \gksolhead/\gkanswer 提取，
% 并与 1980/1980_数学真题_文_答案要点.md 逐题比对后生成，勿手改（改源 MD 后重跑）。

\begin{gknotice}
\textbf{使用说明}：本节只给最终答案便于快速核对，完整推导见解析版。答案为编写者独立推导，\textbf{非官方标准答案}。
\end{gknotice}

\gkpart{题目}
\begin{itemize}[leftmargin=2.6em]
  \item[\textbf{1.}] $\dfrac{9}{13}-\dfrac{7}{13}\mathrm i$
  \item[\textbf{2.}] $x=1$，$y=-2$，$z=3$
  \item[\textbf{3.}] 见解析版
  \item[\textbf{4.}] $32\%$
  \item[\textbf{5.}] $-1$
  \item[\textbf{6.}] (1) 见解析版；(2) 逆命题为“若四边形 $ABCD$ 的对角线 $AC$ 平分对角线 $BD$，则 $AC$ 将四边形分成面积相等的两个三角形”，该逆命题正确
  \item[\textbf{7.}] (1) 见解析版；(2) $AE=6$
  \item[\textbf{8.}] (1) $x^2-\dfrac{y^2}{4}=1$（双曲线，中心在原点，实轴在 $x$ 轴，$a=1$、$b=2$，渐近线 $y=\pm2x$）；(2) 当 $0\le t<\dfrac{\pi}{2}$ 时得到双曲线右支在 $x$ 轴上方的部分（含顶点 $(1,0)$）；当 $\pi\le t<\dfrac{3\pi}{2}$ 时得到左支在 $x$ 轴上方的部分（含顶点 $(-1,0)$）
\end{itemize}
